\chapter{Introduzione}

\noindent La richiesta consiste nel creare un sistema dove gli studenti possono caricare, modificare e
versionare appunti per i corsi. 
Principalmente si tratterà di un editor di testo (semplice markdown o text-area) dotato di cronologia
delle versioni (semplificata) e un'organizzazione per corso/argomento.

\section{As-Is}
L'attuale gestione e condivisione degli appunti tra gli studenti universitari avviene prevalentemente in modo frammentato e manuale. Gli appunti sono dispersi su diverse piattaforme (email, chat, cloud personali), rendendo difficile la reperibilità del materiale. Non avendo un sistema centralizzato è altamente probabile disperdere correzzioni o aggiornamenti importanti. In questo contesto la collaborazione risulta molto complicata e limitata all'iniziativa personale.

\section{To-Be}
Il sistema \textbf{Wiki Collaborativo per Appunti} si propone di a risolvere tali problematiche attraverso un sistema centralizzato di collaborazione tra gli studenti. Esso infatti permetterà di:

\begin{itemize}
    \item \textbf{Condividere il materiale:} Gli appunti sono condivisi in un unico sistema alla portata di tutti e a disposizione di tutti.
    \item \textbf{Collaborare:} Ogni utente, dopo la registrazione, può contribuire aggiugnendo nuovi appunti o modificando gli esistenti.
    \item \textbf{Controllare il materiale:} Ogni modifica agli appunti genera automaticamente una nuova versione. Questo consente agli utenti di navigare nella cronologia e, in caso di errori rilevati, ripristinare la vecchia versione di un appunto.
\end{itemize}